\documentclass[a4paper,12pt,openany]{report}

\usepackage{fontspec}
\setmainfont{XB Niloofar}
\usepackage{fancyhdr}
\usepackage{amsmath}
\usepackage{algorithm}
\usepackage{algorithmic}
\usepackage{graphicx}
\usepackage{hyperref}
\usepackage{xcolor}
\usepackage{titlesec}
\usepackage{setspace}
\usepackage{geometry}
\usepackage{ragged2e} % برای چپ‌چین کردن
\usepackage{lipsum}
\usepackage{polyglossia}
\usepackage{biblatex}
\setdefaultlanguage{farsi}
\setotherlanguage{english}
\usepackage{bidi}
\setRTL 

\geometry{top=3cm, bottom=3cm, left=3cm, right=3cm}

\pagestyle{fancy}
\fancyhf{} %
\fancyhead[LE,RO]{\thepage} 
\fancyhead[LO]{\leftmark} 
\fancyhead[RE]{عنوان پایان‌نامه} 
\fancyfoot[C]{---} 

\onehalfspacing

\setlength{\parskip}{0.5em}

\titleformat{\chapter}[display]
{\normalfont\huge\bfseries}{\chaptertitlename\ \thechapter}{20pt}{\Huge}

\titlespacing*{\chapter}{0pt}{0pt}{40pt}

\let\cleardoublepage\clearpage

\raggedbottom

\addbibresource{references.bib} 

\begin{document}
	
	
	\begin{titlepage}
		\centering
	\includegraphics[width=0.5\textwidth]{logo.png}\par

	\vspace{1cm}
	{\Huge\bfseries {خوشه بندی طیفی با استفاده از بردار ویژه}\par}
	{\Huge\bfseries \LTR{Spectral Clustering with Eigenvector Selection}\par}
	\vspace{0.5cm}
	{\Large استاد راهنما:\par}
	{\large دکتر مهدی دهقان\par}
		\vspace{1cm}
	{\Large نویسنده:\par}
	{\large نیروانا محلو\par}

	
	\vfill
	{\large زمستان 1403\par}
	\end{titlepage}
	
	
	
	\chapter*{چکیده}
	در این پروژه به بررسی و مقایسه روش‌ خوشه‌بندی طیفی  با انتخاب بهینه بردار های ویژه پرداخته ایم .خوشه‌بندی طیفی، به‌عنوان یکی از رویکردهای قدرتمند در تحلیل داده‌ها، امکان شناسایی خوشه‌هایی با ساختارهای پیچیده را فراهم می‌کند. ابتدا مفاهیم اساسی و پیش‌نیازهای این روش، شامل ماتریس‌های شباهت، لاپلاسین و بردارهای ویژه، به تفصیل معرفی شده است. سپس الگوریتم خوشه‌بندی طیفی بر اساس انتخاب بردارهای ویژه مرتبط به‌طور گام‌به‌گام تشریح شده و با الگوریتم‌های k-means و خوشه‌بندی طیفی مبتنی بر کتابخانه مقایسه گردیده است.
	
	
	\tableofcontents
	
	
	\chapter{مقدمه}
	خوشه‌بندی، که به معنای کشف گروه‌بندی‌های طبیعی در داده‌ها است، یکی از جنبه‌های اساسی یادگیری ماشین و تحلیل الگو محسوب می‌شود. این روش به‌ویژه برای تحلیل مجموعه داده‌های بزرگ اهمیت دارد، زیرا به عنوان یک ابزار یادگیری بدون نظارت  نیازی به برچسب‌گذاری دستی و وقت‌گیر داده‌ها ندارد، که اغلب با ناسازگاری همراه است. از جمله روش‌های محبوب خوشه‌بندی می‌توان به مدل‌های ترکیبی (Mixture-Models) و K-Means اشاره کرد که بر اساس مدل‌سازی صریح توزیع داده‌ها عمل می‌کنند. با این حال، این روش‌ها در شرایطی که توزیع داده‌ها ناشناخته، پیچیده یا همراه با نویز باشد، عملکرد ضعیفی دارند.
	
	خوشه‌بندی طیفی به عنوان جایگزینی قدرتمند برای چنین شرایطی معرفی شده است. این روش با استفاده از بردارهای ویژه ماتریس شباهت یا مجاورت که از داده‌ها استخراج می‌شود، خوشه‌بندی را انجام می‌دهد. خوشه‌بندی طیفی حتی در مواردی که روش‌های سنتی مانند مدل‌های ترکیبی یا K-Means قابل اعمال نیستند، عملکرد موفقی دارد. برای مثال، زمانی که داده‌ها توسط بردارهای ویژگی با طول‌های متغیر نمایش داده می‌شوند یا مدل‌سازی صریح توزیع داده‌ها ممکن نیست، خوشه‌بندی طیفی می‌تواند با استفاده از معیارهای شباهت زوجی، داده‌ها را به‌خوبی خوشه‌بندی کند.
	\chapter{مفاهیم کلیدی و پیش‌نیازها}
	\section{تعاریف مفاهیم کلیدی}
	\begin{description}
	\item[\textbullet] \textbf{ماتریس وقوع:}
		ماتریس وقوع برای یک گراف، یک ماتریس \( (0,1) \) است که هر سطر آن نمایانگر یک رأس و هر ستون آن نمایانگر یک یال است. در این ماتریس، مقدار \( (v, e) = 1 \) است اگر و تنها اگر رأس \( v \) بر یال \( e \) واقع باشد.
		.
		\item[\textbullet\ بردار وقوع:] 
		در ریاضیات، \textbf{بردار نشانگر} ، \textbf{بردار مشخصه}  یا \textbf{بردار وقوع}  برای یک زیرمجموعه \( T \) از مجموعه \( S \)، برداری است به صورت 
		\[
		\mathbf{x}_T := (x_s)_{s \in S}
		\]
		که در آن:
		\[
		x_s = 1 \quad \text{، } s \in T, \quad \text{و} \quad x_s = 0 \quad \text{، } s \notin T.
		\]
		
		اگر \( S \) شماراست و عناصر آن شماره‌گذاری شده باشند به صورتی که 
		\( S = \{s_1, s_2, \ldots, s_n\} \)، 
		در این صورت:
		\[
		\mathbf{x}_T = (x_1, x_2, \ldots, x_n),
		\]
		که در آن:
		\[
		x_i = 1 \quad \text{، } s_i \in T, \quad \text{و} \quad x_i = 0 \quad \text{، } s_i \notin T.
		\]
		
		به طور ساده، بردار نشانگر \( T \) برداری است که برای هر عنصر \( S \) یک مقدار دارد، که اگر عنصر موردنظر از \( S \) در \( T \) باشد مقدار آن برابر با یک، و در غیر این صورت صفر است.
		
		
		\vspace{0.5cm}
		
		\item[\textbullet\ کات نرمالایز شده :] 
		 یک روش در علم پردازش تصویر و داده‌کاوی برای تقسیم‌بندی  تصاویر یا گراف‌ها است. این روش یکی از تکنیک‌های پرکاربرد در حوزه‌هایی مثل بینایی کامپیوتر، شبکه‌های اجتماعی و تحلیل داده‌ها به حساب می آید.
		
		
		کات نرمالایز شده سعی می‌کند یک گراف را به چند قسمت تقسیم کند به طوری که:
		
		لبه‌های قطع‌شده (Cut) یا مجموع وزن یال‌های جداکننده کمترین مقدار ممکن باشد و
		تقسیم‌بندی متعادل باشد(یعنی اندازه خوشه‌ها نسبتاً برابر باشد و فقط بخش کوچکی از گراف از دست برود).
		
		مسئله برش نرمالایز شده را می‌توان با یک گراف وزن‌دار \( G = (V, E) \) مدل کرد. وزن \( w_{ij} \) نشان‌دهنده شباهت بین گره‌های \( i \) و \( j \) است. هدف، تقسیم گراف به بخش‌هایی است که گره‌های هر بخش شباهت زیادی داشته باشند و ارتباط کمی با سایر بخش‌ها داشته باشند.
		
		
		
		فرمول برش نرمالایز شده به صورت زیر است:
		\[
		\text{ncut}(A, B) = \frac{w(A, B)}{w(A, V)} + \frac{w(A, B)}{w(B, V)}
		\]
		که \( w(A, B) \) وزن یال‌های بین بخش‌های \( A \) و \( B \) است. 
		
		ارتباط نرمالایز شده (\( \text{nassoc} \)) نیز به صورت زیر تعریف می‌شود:
		\[
		\text{nassoc}(A, B) = \frac{w(A, A)}{w(A, V)} + \frac{w(B, B)}{w(B, V)}
		\]
		
		رابطه بین این دو:
		\[
		\text{ncut}(A, B) = 2 - \text{nassoc}(A, B)
		\]
		\\
		
		در روش برش نرمالایز شده، برای هر برش \( (S, \overline{S}) \) در گراف \( G \)، \( \operatorname{ncut}(S, \overline{S}) \) میزان شباهت بین بخش‌های مختلف را اندازه‌گیری می‌کند، و \( \operatorname{nassoc}(S, \overline{S}) \) میزان شباهت کل گره‌ها در یک بخش را محاسبه می‌کند.
		
		از آنجا که:
		\[
		\operatorname{ncut}(S, \overline{S}) = 2 - \operatorname{nassoc}(S, \overline{S}),
		\]
		یک برش \( (S^*, \overline{S}^*) \) که \( \operatorname{ncut}(S, \overline{S}) \) را کمینه کند، به‌طور همزمان \( \operatorname{nassoc}(S, \overline{S}) \) را بیشینه می‌کند.
		
		محاسبه یک برش \( (S^*, \overline{S}^*) \) که \( \operatorname{ncut}(S, \overline{S}) \) را کمینه کند، یک مسئله \textbf{NP-hard} است. با این حال، می‌توان با استفاده از تکنیک‌های طیفی، در زمان چندجمله‌ای، برشی با وزن نرمالایز کوچک برای \( \operatorname{ncut}(S, \overline{S}) \) پیدا کرد.
		
		الگوریتم Ncut به شکل زیر است.
		
		 داریم:
		\[
		d(i) = \sum_j w_{ij}
		\]
		و ماتریس‌های \( D \) و \( W \) به ترتیب ماتریس درجه و ماتریس وزن هستند. برش نرمالایز شده به صورت زیر بازنویسی می‌شود:
		\[
		\min \frac{y^T (D - W) y}{y^T D y}
		\]
		که \( y \) برداری است که بخش‌های گراف را مشخص می‌کند. 
		
		با ساده‌سازی محدودیت‌ها، مسئله به حل مقدار ویژه تعمیم‌یافته زیر تبدیل می‌شود:
		\[
		(D - W) y = \lambda D y
		\]
		که دومین مقدار ویژه کوچک \( \lambda \) و بردار ویژه متناظر آن، تقسیم‌بندی بهینه را ارائه می‌دهند.
		
		\item[\textbullet\ خوشه بندی پویا:] خوشه‌بندی پویا  روشی در تحلیل داده‌ها و یادگیری ماشین است که شامل گروه‌بندی نقاط داده به خوشه‌هایی است که می‌توانند بر اساس زمان یا ورودی داده‌ها تغییر کنند.
		
		\item[\textbullet\ گراف پویا:] گراف پویا  به گرافی گفته می‌شود که با گذر زمان تغییر می‌کند. این تغییرات می‌توانند شامل اضافه شدن یا حذف گره‌ها ، یال‌ها ، یا تغییر در وزن یال‌ها باشند. گراف‌های پویا برای مدل‌سازی سیستم‌هایی که وضعیت‌شان به مرور زمان تغییر می‌کند، به کار می‌روند.
		\item[\textbullet\ ماتریس شباهت:] 
		ماتریس شباهت  یک ماتریس مربعی و متقارن است که برای نمایش شباهت‌های دوتایی بین اشیاء در یک مجموعه داده استفاده می‌شود. هر عنصر از این ماتریس که با \( A(i, j) \) نمایش داده می‌شود، نشان‌دهنده میزان شباهت بین نقطه داده \( i \) و نقطه داده \( j \) است. عناصر قطر اصلی ماتریس معمولاً بالاترین مقادیر را دارند، زیرا آن‌ها نشان‌دهنده شباهت هر نقطه با خودش هستند که معمولاً بیشینه شباهت ممکن است.
		
		تعریف دقیق "شباهت" بسته به زمینه متفاوت است و می‌تواند با استفاده از معیارهای مختلفی مانند فاصله اقلیدسی، شباهت کسینوسی، همبستگی یا سایر معیارهای خاص حوزه تعریف شود. انتخاب معیار شباهت به طور مستقیم بر ساخت ماتریس شباهت و در نتیجه بر عملکرد الگوریتم یادگیری ماشینی که از آن استفاده می‌کند، تأثیر می‌گذارد.
		
		
		
		ساخت ماتریس شباهت شامل مراحل زیر است:
		\begin{enumerate}
			\item انتخاب یک معیار شباهت مناسب بر اساس ماهیت داده‌ها و مسئله مورد نظر.
			\item محاسبه شباهت‌های دوتایی بین تمام نقاط داده در مجموعه داده با استفاده از معیار انتخاب‌شده.
			\item پر کردن ماتریس با این مقادیر شباهت، به گونه‌ای که ماتریس متقارن باشد، یعنی \( A(i, j) = A(j, i) \).
			\item در صورت نیاز، نرمال‌سازی ماتریس بر اساس الزامات الگوریتم.
		\end{enumerate}
		
		\item[\textbullet\ شباهت گوسی:] 
		
		شباهت گوسی بین دو نقطه داده \( x_i \) و \( x_j \) به صورت زیر تعریف می‌شود:
		\[
		S_{ij} = \exp\left(-\frac{\|x_i - x_j\|^2}{2\sigma^2}\right)
		\]
		
		
		\begin{enumerate}
			\item \( \|x_i - x_j\|^2 \): فاصله اقلیدسی بین نقاط.
			\item \( \sigma \): پارامتر تنظیم که کنترل‌کننده تأثیر فاصله است.
			\item \( \exp(-x) \): تابع نمایی که مقدار شباهت را به بازه \([0, 1]\) محدود می‌کند.
		\end{enumerate}
		
		\item[\textbullet\ گراف شباهت:] گراف شباهت  نموداری از گره‌ها و یال‌ها است که در آن گره‌ها بر اساس میزان شباهتشان چیدمان شده‌اند. هرچه دو گره به یکدیگر شبیه‌تر باشند، در نمودار به هم نزدیک‌تر نمایش داده می‌شوند. اغلب، شباهت به‌عنوان تابعی از یال‌ها تعریف می‌شود، به این صورت که گره‌های با ارتباط قوی‌تر یکدیگر را جذب می‌کنند.
		\item[\textbullet\ مدل ترکیبی:] مدل ترکیبی یا مدل مخلوط  در آمار، یک مدل ترکیبی، نمایش احتمالی یک جمعیت است که حضور زیرجمعیت‌های مختلف را در آن نشان می‌دهد. بر خلاف مدل‌های سنتی، مدل‌های مخلوط نیازی به تخصیص مشاهدات فردی به یک زیرجمعیت خاص ندارند. در عوض، آنها توزیع احتمال مشاهدات را در کل جمعیت توصیف می کنند. مدل‌های مخلوط امکان استنتاج‌های آماری در مورد ویژگی‌های زیرجمعیت‌ها را تنها بر اساس مشاهدات کل جمعیت، بدون هیچ گونه اطلاعاتی در مورد هویت‌های زیرجمعیتی فراهم می‌کنند
		\item[\textbullet\ مدل مخلوط گوسی:]
		
		مدل مخلوط گوسی (GMM) یک مدل آماری است که داده‌ها را به عنوان ترکیبی از چندین توزیع گوسی مدل‌سازی می‌کند. این مدل فرض می‌کند که داده‌ها از \( K \) توزیع گوسی مستقل تشکیل شده‌اند. 
		
		تابع چگالی احتمال کلی برای داده \( x \) به صورت زیر تعریف می‌شود:
		\[
		P(x) = \sum_{k=1}^{K} \pi_k \cdot \mathcal{N}(x | \mu_k, \Sigma_k)
		\]
		که در آن:
		\begin{itemize}
			\item \( K \): تعداد توزیع‌های گوسی.
			\item \( \pi_k \): وزن مخلوط برای توزیع گوسی \( k \)-ام، به طوری که:
			\[
			\sum_{k=1}^{K} \pi_k = 1, \quad \pi_k \geq 0
			\]
			\item \( \mathcal{N}(x | \mu_k, \Sigma_k) \): تابع چگالی احتمال توزیع گوسی:
			\[
			\mathcal{N}(x | \mu_k, \Sigma_k) = \frac{1}{(2\pi)^{d/2} |\Sigma_k|^{1/2}} \exp\left(-\frac{1}{2}(x - \mu_k)^T \Sigma_k^{-1} (x - \mu_k)\right)
			\]
			که در آن \( \mu_k \) میانگین و \( \Sigma_k \) ماتریس کوواریانس است.
		\end{itemize}
		
	\item[\textbullet\ زنجیره و مدل پنهان مارکوف:]
		
		فرض کنید زنجیره‌ای از متغیرهای تصادفی \( X_1, X_2, \ldots \) داشته باشیم که احتمال تغییر وضعیت از زمان \( t \) به \( t+1 \) مستقل از وضعیت‌های قبلی باشد. این زنجیره را زنجیره مارکوف می‌نامیم. به صورت ریاضی، این خاصیت را می‌توان به شکل زیر بیان کرد:
		\[
		\Pr(X_{t+1} = x | X_1 = x_1, X_2 = x_2, \ldots, X_t = x_t) = \Pr(X_{t+1} = x | X_t = x)
		\]
		
		این خاصیت نشان می‌دهد که وضعیت آینده تنها به وضعیت کنونی وابسته است و به وضعیت‌های قبلی وابسته نیست.
		
		
	
		واضح است که در محاسبه این احتمال شرطی باید:
		\[
		\Pr(X_1 = x_1, \ldots, X_t = x_t) > 0
		\]
		باشد. در غیر این صورت، امکان محاسبه احتمال شرطی وجود ندارد. این احتمال به معنای آن است که مسیر \( (X_t \to X_{t+1}) \) با احتمال مثبت در فرآیند تصادفی امکان‌پذیر باشد.
		
		
		زنجیره مارکوف را می‌توان با استفاده از ماتریس انتقال نمایش داد. این ماتریس، احتمالات انتقال از هر وضعیت به وضعیت دیگر را مشخص می‌کند. عناصر این ماتریس به صورت زیر تعریف می‌شوند:
		\[
		P_{ij} = \Pr(X_{t+1} = j | X_t = i)
		\]
		
		در بسیاری از مسائل، زنجیره مارکوف معمولاً از گراف جهت‌دار استفاده می‌کند. مقدار هر یال نمایانگر احتمال انتقال از یک وضعیت به وضعیت دیگر است.
		

		فرض کنید یک زنجیره مارکوف شامل سه وضعیت داشته باشیم. ماتریس انتقال به صورت زیر نوشته می‌شود:
		\[
		P = \begin{bmatrix}
			P_{11} & P_{12} & P_{13} \\
			P_{21} & P_{22} & P_{23} \\
			P_{31} & P_{32} & P_{33}
		\end{bmatrix}
		\]
	
	که در آن \( P_{ij} \) احتمال انتقال از وضعیت \( i \) به \( j \) است.
	
	مدل پنهان مارکوف، درست به مانند یک فرآیند مارکوف است با این تفاوت که فضای حالت در این جا نامعلوم است. به این ترتیب نمی‌توانیم تشخیص دهیم که وضعیت یا حالت با پیشامدها مطابقت دارد. ولی هر وضعیت با یک خروجی همراه است. به این ترتیب قرار است براساس خروجی حدس بزنیم که دنباله مقدارهای فرآیند (وضعیت‌ها) چه بوده است. چنین مدلی بوسیله «لئونارد بام» ریاضی‌دان آمریکایی در سال‌های حدود 1960 تحقیق و مورد بررسی قرار گرفت.
		
	\end{description}
	
		\section{الگوریتم های پیش نیاز}
	   \subsection{K-means}
	   
	 
	   
	   یادگیری ماشین بدون نظارت فرآیند آموزش کامپیوتر برای کار با داده‌های بدون برچسب و دسته‌بندی‌نشده است. در این روش، الگوریتم بدون نظارت عمل کرده و داده‌های نامرتب را بر اساس شباهت‌ها، الگوها و تفاوت‌ها سازماندهی می‌کند. 
	   
	   در روش خوشه‌بندی k-means نقاط داده به یکی از \( K \) خوشه‌ها تخصیص داده می‌شوند، بسته به فاصله آن‌ها از مرکز خوشه‌ها. این الگوریتم با تخصیص تصادفی مراکز خوشه‌ها در فضا شروع می‌شود. سپس هر نقطه داده بر اساس فاصله‌اش از مرکز خوشه به یکی از خوشه‌ها اختصاص داده می‌شود. پس از تخصیص داده‌ها به خوشه‌ها، مراکز خوشه جدید محاسبه و جایگزین می‌شوند. این فرآیند به صورت تکراری ادامه پیدا می‌کند تا خوشه‌بندی بهینه پیدا شود.
	   
	   در تحلیل، معمولاً فرض می‌شود که تعداد خوشه‌ها (\( K \)) از پیش تعیین شده است و داده‌ها باید به یکی از این گروه‌ها تخصیص داده شوند. 
	   
	   در برخی موارد، مقدار \( K \) مشخص نیست و باید مقدار بهینه آن تعیین شود. روش k-means زمانی بهترین عملکرد را دارد که داده‌ها به خوبی از هم جدا باشند. وقتی داده‌ها هم‌پوشانی زیادی دارند، این روش مناسب نیست. k-means در مقایسه با سایر روش‌های خوشه‌بندی سریع‌تر است و اتصال قوی‌تری بین نقاط داده ارائه می‌دهد. با این حال، این روش اطلاعات دقیقی در مورد کیفیت خوشه‌ها ارائه نمی‌دهد. تخصیص اولیه مراکز خوشه می‌تواند به خوشه‌های متفاوتی منجر شود و الگوریتم به نویز حساس است و ممکن است در کمینه محلی گیر کند.
	   
	   \subsubsection*{هدف خوشه‌بندی  K-Means}
	   
	   هدف خوشه‌بندی، تقسیم مجموعه داده‌ها یا جمعیت به چندین گروه است، به طوری که نقاط داده داخل هر گروه شباهت بیشتری به یکدیگر داشته باشند و از نقاط داده در سایر گروه‌ها متفاوت باشند. این اساساً یک گروه‌بندی بر اساس شباهت‌ها و تفاوت‌ها است.
	   
	   \subsubsection*{نحوه کار الگوریتم  K-Means}
	   
	   فرض کنید یک مجموعه داده شامل آیتم‌هایی با ویژگی‌ها و مقادیر مشخص داریم (مثل یک بردار). وظیفه این است که این آیتم‌ها را به گروه‌هایی دسته‌بندی کنیم. برای این کار، از الگوریتم K-means یک الگوریتم یادگیری بدون نظارت، استفاده می‌کنیم. 
	   
	   \( K \) در نام الگوریتم، تعداد گروه‌ها یا خوشه‌هایی را که می‌خواهیم داده‌ها در آن‌ها دسته‌بندی شوند، نشان می‌دهد.
	   
	   (می‌توانید آیتم‌ها را به صورت نقاط در یک فضای \( n \)-بعدی تصور کنید). الگوریتم آیتم‌ها را به \( K \) گروه یا خوشه براساس شباهت آن‌ها دسته‌بندی می‌کند. برای محاسبه شباهت، از فاصله اقلیدسی استفاده می‌کنیم.
	   
	   \subsubsection*{مراحل الگوریتم}
	   
	   الگوریتم به صورت زیر عمل می‌کند:
	   \begin{enumerate}
	   	\item ابتدا \( K \) نقطه به صورت تصادفی به عنوان مراکز خوشه (centroids) مقداردهی اولیه می‌شوند.
	   	\item هر آیتم به نزدیک‌ترین مرکز خوشه تخصیص داده می‌شود. مختصات مراکز خوشه با میانگین نقاط اختصاص داده شده به آن خوشه به‌روزرسانی می‌شود.
	   	\item این فرآیند برای تعداد مشخصی از تکرارها ادامه می‌یابد. در نهایت، خوشه‌ها تشکیل می‌شوند.
	   \end{enumerate}
	   
	   نقاطی که به عنوان مراکز خوشه شناخته می‌شوند، میانگین مقادیر آیتم‌های تخصیص‌یافته به آن‌ها هستند. برای مقداردهی اولیه این مراکز، روش‌های مختلفی وجود دارد. یکی از روش‌های ساده این است که مراکز را از بین آیتم‌های موجود به صورت تصادفی انتخاب کنیم. روش دیگر این است که مقادیر اولیه را به صورت تصادفی در بازه مقادیر داده‌ها تعیین کنیم (برای مثال، اگر ویژگی \( x \) مقادیری در بازه \([0, 3]\) داشته باشد، مراکز خوشه نیز می‌توانند در همین بازه مقداردهی شوند).
	   
	 \subsection{Expectation-Maximization}
	 
	 در برنامه‌های واقعی یادگیری ماشین، معمول است که بسیاری از ویژگی‌ها مهم باشند، اما تنها زیرمجموعه‌ای از آن‌ها قابل مشاهده باشد. وقتی با متغیرهایی مواجه می‌شویم که گاهی قابل مشاهده هستند و گاهی نه، می‌توان از مواردی که متغیر مشاهده می‌شود برای یادگیری الگوها و پیش‌بینی مقادیر آن در مواردی که غیرقابل مشاهده است، استفاده کرد. این رویکرد معمولاً به‌عنوان «مدیریت داده‌های گمشده» شناخته می‌شود. الگوریتم بیشینه‌سازی انتظار (EM) یکی از روش‌های اصلی برای حل این مسئله است.
	 
	 الگوریتم EM برای متغیرهای پنهان کاربرد دارد. این متغیرها قابل مشاهده مستقیم نیستند، اما از مقادیر متغیرهای مشاهده‌شده استنتاج می‌شوند. EM یک چارچوب برای یافتن پارامترهای بیشینه درست‌نمایی در مدل‌های آماری ارائه می‌دهد.
	 
	 \subsubsection*{متغیرهای پنهان و مفاهیم کلیدی}
	 
	 \begin{itemize}
	 	\item \textbf{متغیرهای پنهان:} متغیرهایی که به طور مستقیم قابل مشاهده نیستند، اما از اثر آن‌ها بر متغیرهای مشاهده‌شده می‌توان اطلاعات کسب کرد.
	 	\item \textbf{درست‌نمایی:} احتمال مشاهده داده‌ها بر اساس پارامترهای مدل. هدف الگوریتم EM یافتن پارامترهایی است که درست‌نمایی را بیشینه کنند.
	 	\item \textbf{لگاریتم درست‌نمایی:} لگاریتم تابع درست‌نمایی که تطابق بین داده‌های مشاهده‌شده و مدل را اندازه‌گیری می‌کند.
	 	\item \textbf{بیشینه‌سازی درست‌نمایی (MLE):} روشی برای یافتن پارامترهایی که تابع درست‌نمایی را بیشینه می‌کنند.
	 	\item \textbf{احتمال پسین:} در بافت استنتاج بیزی، EM می‌تواند برای تخمین بیشینه احتمال پسین (MAP) توسعه یابد.
	 \end{itemize}
	 
	 الگوریتم EM از دو مرحله اصلی تشکیل شده است که به صورت تکراری اجرا می‌شوند:
	 
	 \subsubsection*{مرحله انتظار (E-Step)}
	 در این مرحله، مقادیر مورد انتظار یا احتمال پسین متغیرهای پنهان با استفاده از داده‌های مشاهده‌شده و مقادیر فعلی پارامترها محاسبه می‌شود:
	 \begin{itemize}
	 	\item احتمال پسین یا مسئولیت هر متغیر پنهان برای هر نقطه داده محاسبه می‌شود.
	 	\item مقادیر گمشده یا ناقص داده‌ها تخمین زده می‌شوند.
	 	\item لگاریتم درست‌نمایی داده‌های مشاهده‌شده بر اساس مقادیر فعلی پارامترها محاسبه می‌شود.
	 \end{itemize}
	 
	 \subsubsection*{مرحله بیشینه‌سازی (M-Step)}
	 در این مرحله، پارامترهای مدل به‌روزرسانی می‌شوند:
	 \begin{itemize}
	 	\item پارامترهای مدل با بیشینه کردن لگاریتم درست‌نمایی کامل داده‌ها به‌روزرسانی می‌شوند.
	 	\item این مرحله معمولاً شامل حل مسائل بهینه‌سازی برای یافتن مقادیر پارامترهایی است که لگاریتم درست‌نمایی را بیشینه می‌کنند.
	 \end{itemize}
	 
	 \subsubsection*{همگرایی (Convergence)}
	 در این مرحله بررسی می‌شود که آیا الگوریتم به یک راه‌حل پایدار رسیده است یا خیر:
	 \begin{itemize}
	 	\item اگر تغییر در لگاریتم درست‌نمایی یا مقادیر پارامترها کمتر از یک مقدار آستانه از پیش تعریف‌شده باشد، الگوریتم متوقف می‌شود.
	 	\item در غیر این صورت، به مرحله \( E \) بازگشته و فرآیند تا رسیدن به همگرایی ادامه پیدا می‌کند.
	 \end{itemize}
	 
	 \subsection*{کاربردهای الگوریتم EM}
	 
	 \begin{itemize}
	 	\item خوشه‌بندی داده‌ها (Clustering)
	 	\item مدل‌سازی آماری
	 	\item تشخیص الگو در بینایی کامپیوتر
	 	\item پردازش زبان طبیعی
	 \end{itemize}
	 
	 در ادامه الگوریتم EM برای خوشه‌بندی داده‌ها با استفاده از مدل مخلوط گوسی (GMM) توضیح داده می‌شود. این الگوریتم از دو مرحله اصلی \textbf{انتظار} (E-Step) و \textbf{بیشینه‌سازی} (M-Step) برای شناسایی پارامترهای مدل استفاده می‌کند.
	 
	 \subsubsection*{مراحل الگوریتم}
	
	 	۱. \textbf{ورودی:} مجموعه داده $\{x_1, x_2, \dots, x_N\}$ و تعداد خوشه‌ها $K$.
	 	
	 	۲. \textbf{مقداردهی اولیه:} مقادیر اولیه $\mu_k, \Sigma_k, w_k$ برای $k=1, \dots, K$ به صورت تصادفی تعیین می‌شوند.
	 	
	 	۳. تکرار تا رسیدن به همگرایی:
	 	\begin{itemize}
	 		\item \textbf{مرحله انتظار (E-Step):}
	 		\begin{itemize}
	 			\item برای هر نقطه داده $x_i$ و هر خوشه $k$، احتمال مسئولیت محاسبه می‌شود:
	 			\[
	 			P(z_k | x_i) = \frac{w_k \cdot \mathcal{N}(x_i | \mu_k, \Sigma_k)}{\sum_{j=1}^{K} w_j \cdot \mathcal{N}(x_i | \mu_j, \Sigma_j)}
	 			\]
	 		\end{itemize}
	 		\item \textbf{مرحله بیشینه‌سازی (M-Step):}
	 		\begin{itemize}
	 			\item به‌روزرسانی میانگین خوشه:
	 			\[
	 			\mu_k = \frac{\sum_{i=1}^{N} P(z_k | x_i) \cdot x_i}{\sum_{i=1}^{N} P(z_k | x_i)}
	 			\]
	 			\item به‌روزرسانی ماتریس کواریانس:
	 			\[
	 			\Sigma_k = \frac{\sum_{i=1}^{N} P(z_k | x_i) \cdot (x_i - \mu_k)(x_i - \mu_k)^T}{\sum_{i=1}^{N} P(z_k | x_i)}
	 			\]
	 			\item به‌روزرسانی وزن خوشه:
	 			\[
	 			w_k = \frac{\sum_{i=1}^{N} P(z_k | x_i)}{N}
	 			\]
	 		\end{itemize}
	 	\end{itemize}
	 	
	 	۴. \textbf{خروجی:} پارامترهای $\mu_k, \Sigma_k, w_k$ و تخصیص داده‌ها به خوشه‌ها.

\chapter{ماتریس لاپلاسی}
\section*{گراف های ساده}


فرض کنید یک گراف ساده \( G \) با \( n \) راس \( v_1, v_2, \dots, v_n \) داده شده است. ماتریس لاپلاسین \( L_{n \times n} \) به صورت درایه‌ای به شکل زیر تعریف می‌شود:

\[
L_{i,j} := \begin{cases} 
	\deg(v_i) &  i = j, \\
	-1 & \text i \neq j \text{ ، } v_i \text{ به } v_j \text{ متصل باشد}, \\
	0 & \text{در غیر این صورت.}
\end{cases}
\]

یا به صورت معادل توسط ماتریس زیر:

\[
L = D - A,
\]

که در آن:
\begin{itemize}
	\item \( D \) ماتریس درجه  است.
	\item \( A \) ماتریس همسایگی  گراف است.
\end{itemize}

از آنجا که \( G \) یک گراف ساده است، \( A \) تنها شامل مقادیر \( 0 \) و \( 1 \) است و عناصر قطری آن همگی برابر با \( 0 \) هستند.
\\
	
\includegraphics[width=1\textwidth]{1.png}\par\vspace{1cm}
\includegraphics[width=1\textwidth]{2.png}\par\vspace{1cm}


	
	\section*{ماتریس لاپلاسین برای یک گراف بدون جهت با استفاده از ماتریس وقوع جهت‌دار}
	
	ماتریس وقوع جهت‌دار \( B \) با ابعاد \( |v| \times |e| \) و درایه \( B_{ve} \) برای رأس \( v \) و یال \( e \) (که رأس‌های \( v_i \) و \( v_j \) را متصل می‌کند، با \( i \neq j \)) به صورت زیر تعریف می‌شود:
	
	\[
	B_{ve} = \begin{cases} 
		1, &  v = v_i, \\
		-1, &  v = v_j, \\
		0, & \text{در غیر این صورت.}
	\end{cases}
	\]
	
	با اینکه یال‌ها در این تعریف به صورت فنی جهت‌دار هستند، جهت آن‌ها می‌تواند دلخواه باشد و همچنان به همان ماتریس لاپلاسین متقارن \( |v| \times |v| \) منجر می‌شود که به صورت زیر تعریف می‌شود:
	
	\[
	L = B B^T
	\]
	
	که در آن \( B^T \) ترانهاده ماتریس \( B \) است.
	
\includegraphics[width=1\textwidth]{3.png}\par\vspace{1cm}

\section*{لاپلاسین متقارن برای گراف جهت‌دار}

ماتریس لاپلاسین یک گراف جهت‌دار به‌صورت پیش‌فرض معمولاً نامتقارن است، در حالی که تکنیک‌هایی مانند خوشه‌بندی طیفی معمولاً برای گراف‌های بدون جهت با ماتریس‌های همسایگی و لاپلاسین متقارن طراحی شده‌اند. یک روش ساده برای استفاده از این تکنیک‌ها، تبدیل گراف جهت‌دار اولیه به یک گراف بدون جهت و ساخت ماتریس لاپلاسین برای گراف تبدیل‌شده است.

در نمایش ماتریسی، ماتریس همسایگی گراف بدون جهت می‌تواند به صورت یک جمع بولی از ماتریس همسایگی \( A \) گراف جهت‌دار اولیه و ترانهاده آن \( A^T \) تعریف شود، که در آن عناصر \( 0 \) و \( 1 \) ماتریس \( A \) به عنوان مقادیر منطقی (به جای مقادیر عددی) در نظر گرفته می‌شوند. به صورت مثال زیر:

\[
A_{\text{undirected}} = A \lor A^T
\]

که \( \lor \) نشان‌دهنده عملیات OR منطقی است. این تبدیل باعث می‌شود که تقارن لازم برای ساخت ماتریس لاپلاسین به دست آید.
\\

\includegraphics[width=1\textwidth]{4.png}\par\vspace{1cm}

	
	\section*{نرمال‌سازی ماتریس لاپلاسین}
	
	رأسی با درجه زیاد، که گاهی به آن رأس سنگین (Heavy-Node) نیز گفته می‌شود، باعث می‌شود که ورودی‌های قطری بزرگی در ماتریس لاپلاسین ایجاد شود و ویژگی‌های ماتریس را تحت تأثیر قرار دهد. هدف از نرمال‌سازی، برابرتر کردن تأثیر چنین رأس‌هایی با سایر رأس‌ها است. این کار با تقسیم درایه‌های ماتریس لاپلاسین بر درجه رأس‌ها انجام می‌شود. برای جلوگیری از تقسیم بر صفر، رأس‌های ایزوله با درجه صفر از فرآیند نرمال‌سازی حذف می‌شوند.
	
	\subsection*{لاپلاسین نرمال‌شده متقارن}
	
	ماتریس لاپلاسین نرمال‌شده متقارن به صورت زیر تعریف می‌شود:
	\[
	L^{\text{sym}} := (D^+)^{1/2} L (D^+)^{1/2} = I - (D^+)^{1/2} A (D^+)^{1/2},
	\]
	
	که در آن \( D^+ \) معکوس Moore-Penrose ماتریس درجه است.
	
	درایه‌های \( L^{\text{sym}} \) به صورت زیر داده می‌شوند:
	\[
	L^{\text{sym}}_{i,j} := \begin{cases} 
		1, &  i = j \text{ و } \deg(v_i) \neq 0, \\
		-\frac{1}{\sqrt{\deg(v_i) \deg(v_j)}}, &  i \neq j \text{ و } v_i \text{ به } v_j \text{ متصل باشد}, \\
		0, & \text{در غیر این صورت.}
	\end{cases}
	\]
	
	ماتریس لاپلاسین نرمال‌شده متقارن تنها در صورتی متقارن است که ماتریس همسایگی متقارن باشد.
\includegraphics[width=1\textwidth]{5.png}\par\vspace{1cm}
\includegraphics[width=1\textwidth]{6.png}\par\vspace{1cm}
\section*{لاپلاسین نرمال‌شده چپ (تصادفی) و راست}

\subsection*{لاپلاسین نرمال‌شده چپ}
لاپلاسین نرمال‌شده چپ (Random-Walk) به صورت زیر تعریف می‌شود:
\[
L^{\text{rw}} := D^+ L = I - D^+ A,
\]
که در آن \( D^+ \) معکوس Moore-Penrose ماتریس درجه است. درایه‌های \( L^{\text{rw}} \) به صورت زیر داده می‌شوند:
\[
L^{\text{rw}}_{i,j} := \begin{cases} 
	1, & \text i = j \text{ و } \deg(v_i) \neq 0, \\
	-\frac{1}{\deg(v_i)}, & \text i \neq j \text{ و } v_i \text{ به } v_j \text{ متصل باشد}, \\
	0, & \text{در غیر این صورت.}
\end{cases}
\]

\subsection*{لاپلاسین نرمال‌شده راست}
به طور مشابه، لاپلاسین نرمال‌شده راست به صورت زیر تعریف می‌شود:
\[
L^{\text{right}} = LD^+ = I - AD^+.
\]

\subsection*{ویژگی‌ها}
- لاپلاسین نرمال‌شده چپ (\( L^{\text{rw}} \)) و راست (\( L^{\text{right}} \)) معمولاً در صورتی که ماتریس همسایگی متقارن نباشد، متقارن نیستند، مگر در حالت خاصی که تمامی رأس‌ها ایزوله باشند.
- اگر گراف \( G \) هیچ رأسی ایزوله نداشته باشد، \( D^+ A \) یک ماتریس تصادفی راست (Right-Stochastic) است و در نتیجه \( L^{\text{rw}} \) دارای مجموع سطری صفر است.
- به طور مشابه، \( AD^+ \) یک ماتریس تصادفی چپ (Left-Stochastic) است و \( L^{\text{right}} \) دارای مجموع ستونی صفر است.

\subsection*{رابطه با ماتریس‌های تصادفی}
- لاپلاسین نرمال‌شده چپ مرتبط با ماتریس تصادفی راست \( D^+_\text{out} A \) است، که مجموع سطری صفر دارد.
- لاپلاسین نرمال‌شده راست مرتبط با ماتریس تصادفی چپ \( AD^+_\text{in} \) است، که مجموع ستونی صفر دارد.
\includegraphics[width=1\textwidth]{7.png}\par\vspace{1cm}
\includegraphics[width=1\textwidth]{8.png}\par\vspace{1cm}

\section*{ گراف‌ها با یال‌های وزن‌دار}

در کاربردهای معمول، گراف‌هایی با یال‌های وزن‌دار به‌سادگی توسط ماتریس‌های همسایگی آن‌ها تعریف می‌شوند، که مقادیر درایه‌ها عددی بوده و دیگر محدود به مقادیر صفر و یک نیستند. در روش‌هایی مانند خوشه‌بندی طیفی  و پردازش سیگنال مبتنی بر گراف  که رأس‌های گراف نشان‌دهنده نقاط داده هستند، وزن یال‌ها می‌تواند به‌صورت معکوس متناسب با فاصله بین جفت نقاط داده محاسبه شود. این کار منجر به وزن‌های غیرمنفی می‌شود، به‌طوری‌که مقادیر بزرگ‌تر به‌طور غیررسمی به جفت نقاط داده مشابه‌تر اشاره دارند.

استفاده از همبستگی و ضدهمبستگی بین نقاط داده به‌طور طبیعی به وزن‌های مثبت و منفی منجر می‌شود. بیشتر تعاریف برای گراف‌های ساده به‌راحتی به حالت استاندارد وزن‌های غیرمنفی تعمیم می‌یابند، در حالی که وزن‌های منفی نیاز به توجه بیشتری دارند، به‌ویژه در فرآیند نرمال‌سازی.
\\
\includegraphics[width=1\textwidth]{9.png}\par\vspace{1cm}

\section*{لاپلاسین وزن دار متقارن با استفاده از ماتریس وقوع}

برای گراف‌هایی با یال‌های وزن‌دار، می‌توان یک ماتریس وقوع وزن‌دار \( B \) تعریف کرد و از آن برای ساخت لاپلاسین متقارن متناظر به صورت زیر استفاده کرد:
\[
L = B W B^T.
\]

یک روش جایگزین و تمیزتر، جدا کردن وزن‌ها از اتصالات است: ادامه استفاده از ماتریس وقوع \( B \) برای گراف‌های معمول و معرفی یک ماتریس \( W \) که تنها مقادیر وزن‌ها را نگه می‌دارد. سیستم فنر (Spring System) نمونه‌ای از این مدل است که در مکانیک برای توصیف یک سیستم فنرها با سختی‌ها و طول واحد معین استفاده می‌شود، که در آن مقادیر سختی نقش وزن یال‌های گراف را بازی می‌کنند.

ماتریس وقوع \( B \) بدون وزن برای رأس \( v \) و یال \( e \) (که رأس‌های \( v_i \) و \( v_j \) را متصل می‌کند، با \( i > j \)) به صورت زیر تعریف می‌شود:
\[
B_{ve} = \begin{cases} 
	1, &  v = v_i, \\
	-1, &  v = v_j, \\
	0, & \text{در غیر این صورت.}
\end{cases}
\]

ما همچنین یک ماتریس قطری \( W \) با ابعاد \( |e| \times |e| \) تعریف می‌کنیم که وزن یال‌ها را شامل می‌شود. با اینکه یال‌ها در تعریف \( B \) به‌صورت فنی جهت‌دار هستند، جهت آن‌ها می‌تواند دلخواه باشد و همچنان به همان ماتریس لاپلاسین متقارن \( |v| \times |v| \) منجر می‌شود که به صورت زیر تعریف می‌شود:
\[
L = B W B^T,
\]

که در آن \( B^T \) ترانهاده ماتریس \( B \) است.
\\
\includegraphics[width=1\textwidth]{10.png}\par\vspace{1cm}
\section*{لاپلاسین متقارن برای گراف وزن دار جهت‌دار}

مشابه گراف‌های ساده، ماتریس لاپلاسین یک گراف وزن‌دار جهت‌دار به‌صورت پیش‌فرض معمولاً نامتقارن است. تقارن می‌تواند با تبدیل گراف جهت‌دار اولیه به یک گراف بدون جهت قبل از ساخت لاپلاسین، تضمین شود. ماتریس همسایگی گراف بدون جهت می‌تواند به‌عنوان مجموع ماتریس همسایگی \( A \) گراف جهت‌دار اولیه و ترانهاده آن \( A^T \) تعریف شود:
\[
A_{\text{undirected}} = A + A^T.
\]

در اینجا، عناصر صفر و یک ماتریس \( A \) به‌عنوان مقادیر عددی در نظر گرفته می‌شوند، نه منطقی (Logical) همانند گراف‌های ساده. این تفاوت در نتایج را توضیح می‌دهد: برای گراف‌های ساده، گراف متقارن‌شده همچنان باید ساده باقی بماند و ماتریس همسایگی متقارن آن تنها شامل مقادیر منطقی باشد (مثلاً، جمع منطقی \( 1 \lor 1 = 1 \))، درحالی‌که جمع عددی \( 1 + 1 = 2 \) است.

روش جایگزین برای محاسبه لاپلاسین متقارن استفاده از دو لاپلاسین با درجات ورودی و خروجی است:
\[
L_{\text{sym}} = \frac{1}{2}(L_{\text{in}} + L_{\text{out}}),
\]
که در آن \( L_{\text{in}} \) و \( L_{\text{out}} \) به ترتیب لاپلاسین‌های مرتبط با درجات ورودی و خروجی هستند.
\\
\includegraphics[width=1\textwidth]{11.png}\par\vspace{1cm}

\includegraphics[width=1\textwidth]{12.png}\par\vspace{1cm}

\section*{(وزن دار)نرمال‌سازی ماتریس لاپلاسین}

هدف از نرمال‌سازی، مشابه گراف‌های ساده، این است که مقادیر قطری ماتریس لاپلاسین به واحد تبدیل شوند و در عین حال مقادیر خارج از قطر نیز متناسب تنظیم شوند. در یک گراف وزن‌دار، یک رأس ممکن است درجه بزرگی داشته باشد به دلیل تعداد کمی از یال‌های متصل با وزن‌های بزرگ یا تعداد زیادی از یال‌های متصل با وزن واحد.

حلقه‌های خودی گراف (Self-Loops)، یعنی درایه‌های غیرصفر در قطر اصلی ماتریس همسایگی، بر مقادیر ماتریس لاپلاسین تأثیری ندارند، اما ممکن است نیاز باشد برای محاسبه فاکتورهای نرمال‌سازی لحاظ شوند.

\subsection*{لاپلاسین نرمال‌شده متقارن}

لاپلاسین نرمال‌شده متقارن به صورت زیر تعریف می‌شود:
\[
L^{\text{sym}} := (D^+)^{1/2} L (D^+)^{1/2} = I - (D^+)^{1/2} A (D^+)^{1/2},
\]

که در آن:
\begin{itemize}
	\item \( L \): لاپلاسین غیرنرمال‌شده.
	\item \( A \): ماتریس همسایگی.
	\item \( D \): ماتریس درجه.
	\item \( D^+ \): معکوس Moore-Penrose ماتریس \( D \).
\end{itemize}

از آنجا که ماتریس \( D \) قطری است، ریشه دوم معکوس قطری \( (D^+)^{1/2} \) تنها ماتریس قطری است که درایه‌های قطری آن عکس ریشه‌های دوم درایه‌های قطری \( D \) هستند. اگر تمام وزن‌های یال غیرمنفی باشند، تمامی مقادیر درجه خودکار غیرمنفی خواهند بود و در نتیجه هر مقدار درجه یک ریشه دوم مثبت یکتا خواهد داشت. برای جلوگیری از تقسیم بر صفر، رأس‌هایی با درجه صفر از فرآیند نرمال‌سازی حذف می‌شوند.

لاپلاسین نرمال‌شده متقارن یک ماتریس متقارن است، اگر و تنها اگر ماتریس همسایگی \( A \) متقارن باشد و درایه‌های قطری \( D \) غیرمنفی باشند. در این حالت، می‌توان از اصطلاح لاپلاسین نرمال‌شده متقارن استفاده کرد.

لاپلاسین نرمال‌شده متقارن همچنین می‌تواند به صورت زیر نوشته شود:
\[
L^{\text{sym}} := (D^+)^{1/2} L (D^+)^{1/2} = (D^+)^{1/2} B W B^T (D^+)^{1/2} = S S^T,
\]

که در آن:
\begin{itemize}
	\item \( B \): ماتریس وقوع بدون وزن \( |v| \times |e| \).
	\item \( W \): ماتریس قطری \( |e| \times |e| \) شامل وزن‌های یال‌ها.
	\item \( S \): ماتریس وقوع وزن‌دار جدید \( |v| \times |e| \).
\end{itemize}

ماتریس \( S \) به صورت \( S = (D^+)^{1/2} B W^{1/2} \) تعریف می‌شود که سطرهای آن توسط رأس‌ها و ستون‌های آن توسط یال‌های \( G \) ایندکس می‌شوند.

برای یک یال \( e = \{u, v\} \)، یک درایه \( \frac{1}{\sqrt{d_u}} \) در سطر مربوط به \( u \)، یک درایه \( -\frac{1}{\sqrt{d_v}} \) در سطر مربوط به \( v \) و درایه‌های دیگر صفر هستند.
\\
\includegraphics[width=1\textwidth]{13.png}\par\vspace{1cm}


\section*{(وزن دار)لاپلاسین نرمال‌شده تصادفی (Random-Walk)}

لاپلاسین نرمال‌شده تصادفی به صورت زیر تعریف می‌شود:
\[
L^{\text{rw}} := D^+ L = I - D^+ A,
\]

که در آن:
\begin{itemize}
	\item \( D \): ماتریس درجه 
	\item \( D^+ \): معکوس Moore-Penrose ماتریس \( D \)، که یک ماتریس قطری است و درایه‌های آن معکوس مقادیر قطری \( D \) هستند.
	\item \( A \): ماتریس همسایگی.
	\item \( L \): ماتریس لاپلاسین غیرنرمال‌شده.
\end{itemize}

برای رأس‌های ایزوله (با درجه صفر)، معمولاً مقدار متناظر در \( L^{\text{rw}} \) برابر صفر در نظر گرفته می‌شود. درایه‌های \( L^{\text{rw}} \) به صورت زیر تعریف می‌شوند:
\[
L^{\text{rw}}_{i,j} := \begin{cases} 
	1, & \text i = j \text{ و } \deg(v_i) \neq 0, \\
	-\frac{1}{\deg(v_i)}, & \text i \neq j \text{ و } v_i \text{ به } v_j \text{ متصل باشد}, \\
	0, & \text{در غیر این صورت.}
\end{cases}
\]

نام \( L^{\text{rw}} \) از این واقعیت ناشی می‌شود که این ماتریس به صورت \( L^{\text{rw}} = I - P \) است، که در آن \( P = D^+ A \) به‌سادگی ماتریس انتقال یک Random-Walk روی گراف است که از وزن‌های غیرمنفی تشکیل شده است. به عنوان مثال، فرض کنید \( e_i \) نشان‌دهنده بردار پایه استاندارد \( i \)-ام باشد. در این صورت، \( x = e_i \) یک بردار احتمال است که توزیع مکان‌های Random-Walker را پس از انجام یک گام از رأس \( i \) نشان می‌دهد.

به طور کلی، اگر \( x \) یک توزیع احتمال از مکان Random-Walker روی رأس‌های گراف باشد، آنگاه \( x' = x P \) توزیع احتمال Walker پس از \( t \) گام را نشان می‌دهد.

لاپلاسین نرمال‌شده تصادفی \( L^{\text{rw}} \) همچنین می‌تواند به صورت \( D^+ L \) نوشته شود، زیرا نرمال‌سازی با ضرب \( D^+ \) از سمت چپ انجام می‌شود. هر سطر مجموعی برابر صفر دارد زیرا \( P = D^+ A \) ماتریسی Right-Stochastic است (فرض بر این است که همه وزن‌ها غیرمنفی هستند).

در لاپلاسین نرمال‌شده راست کمتر رایج \( L D^+ = I - A D^+ \)، هر ستون مجموعی برابر صفر دارد زیرا \( A D^+ \) ماتریسی Left-Stochastic است.
\\
\includegraphics[width=1\textwidth]{14.png}\par\vspace{1cm}


	\chapter{خوشه بندی طیفی}

	

مجموعه‌ای از \( N \) نقطه داده داریم که به‌صورت بردار ویژگی‌های زیر تعریف شده‌اند:

\[
D = \{f_1, f_2, \dots, f_N\}.
\]

هر نقطه داده \( f_i \) یک بردار ویژگی است.  
هدف اصلی خوشه‌بندی طیفی، کشف گروه‌بندی طبیعی این داده‌ها است.  
تعداد خوشه‌ها (\( K_0 \)) باید به‌گونه‌ای انتخاب شود که توزیع داده‌ها را به‌بهترین شکل ممکن توصیف کند. اگر مقدار \( K_0 \) درست تخمین زده شود، برابر با تعداد واقعی خوشه‌ها (\( K_{\text{true}} \)) خواهد بود . توجه داشته باشید
که بردارهای ویژگی مختلف می توانند ابعاد متفاوتی داشته باشند.

\section{ساخت ماتریس مشابهت (\( A \))}

ماتریس مشابهت \( A \) به‌صورت \( N \times N \) تعریف می‌شود.  
هر عنصر \( A_{ij} \) نشان‌دهنده میزان شباهت یا وابستگی بین نقاط داده \( i \) و \( j \) است.  
ماتریس \( A \) باید متقارن باشد (\( A_{ij} = A_{ji} \)).  
از بردارهای ویژه (Eigenvectors) ماتریس \( A \) می‌توان برای خوشه‌بندی مستقیم استفاده کرد 
امااستفاده از بردارهای ویژه ماتریس نرمال‌شده مشابهت (\( \tilde{A} \)) بهتر است.
حالا این ماتریس را به شکلی تغییر می‌دهیم که اثر اندازه‌های مختلف نقاط داده متعادل شود. 

\section{ ماتریس نرمال‌شده (\( \tilde{A} \))}

ماتریس مشابهت نرمال‌شده \( \tilde{A} \) (لاپلاسی نرمال شده) به‌صورت زیر تعریف می‌شود:

\[
\tilde{A} = D^{-\frac{1}{2}} A D^{-\frac{1}{2}},
\]

که در آن   \( D\) ماتریس قطری است، یک ماتریس قطری که مقدار قطر اصلی آن برابر است با مجموع وزن‌های یال‌های هر گره:
	\[
	L_{ii} = \sum_j A_{ij}.
	\]


عمل نرمال‌سازی در اینجا برای کاهش اثر مقیاس داده‌ها و بهبود کیفیت خوشه‌بندی انجام می‌شود.

پس از محاسبه ماتریس نرمال‌شده \( \tilde{A} \)، مقادیر ویژه و بردارهای ویژه آن محاسبه می‌شوند.  
\( K_m \) تعداد تخمینی خوشه‌ها است که باید بزرگ‌تر از تعداد واقعی خوشه‌ها (\( K_0 \)) باشد.  
این مقدار برای انتخاب بزرگ‌ترین بردارهای ویژه (بردارهایی که با مقادیر ویژه بزرگ‌تر متناظر هستند) استفاده می‌شود.  
با استفاده از \( K_m \) بردار ویژه بزرگ، داده‌ها به یک فضای جدید با ابعاد کاهش‌یافته منتقل می‌شوند.  

داده‌های انتقال‌یافته به صورت زیر تعریف می‌شوند:
\[
D_e = \{x_1, x_2, \dots, x_N\},
\]
که در آن \( x_n \) یک بردار \( K_m \)-بعدی است.  

هر بردار \( x_n \) به صورت زیر تعریف می‌شود:
\[
x_n = [e_{n1}, e_{n2}, \dots, e_{nK_m}],
\]
که در آن \( e_{nk} \) عنصر \( n \)-ام از \( k \)-امین بردار ویژه است.

اکنون هر بردار ویژگی در مجموعه داده جدید دارای همان ابعاد \( K_m \) است.  
وظیفه خوشه‌بندی طیفی در این مرحله تعیین تعداد خوشه‌ها و سپس گروه‌بندی داده‌ها به خوشه‌های مختلف با استفاده از نمایش جدید داده‌ها در فضای ویژه (eigenspace) است.
\section{بردار ویژه مناسب}
در فضای اصلی (فضای اولیه‌ای که داده‌ها قرار دارند)، ممکن است داده‌هایی که به یک خوشه تعلق دارند، از هم خیلی دور به نظر برسند. برای اینکه کار خوشه‌بندی راحت‌تر شود، داده‌ها را به یک فضای جدید و ساده‌تر منتقل می‌کنیم که در آن نقاط مشابه به هم نزدیک‌تر باشند.

فقط یک زیرمجموعه از \( K_m \) بزرگ‌ترین بردارهای ویژه برای گروه‌بندی \( K_o \) خوشه‌ها مرتبط است.  
مهم است که ابتدا بردارهای ویژه نامربوط یا غیرمفید را شناسایی و حذف کنیم، سپس عملیات خوشه‌بندی را انجام دهیم.  

چگونه می‌توانیم ارتباط یک بردار ویژه را اندازه‌گیری کنیم؟  
یک راه‌حل شهودی این است که مقدار ویژه مربوط به آن بردار ویژه را بررسی کنیم.

در حالت ایده‌آل، زمانی که خوشه‌ها کاملاً از یکدیگر جدا و مشخص هستند، بردارهای ویژه مرتبط که به تعداد خوشه‌ها (\( K_{\text{true}} \)) هستند، مقادیر ویژه‌ای برابر با 1 دارند. سایر بردارهای ویژه مقادیر ویژه کوچک‌تری دارند. در چنین شرایطی، می‌توان به‌سادگی با بررسی مقادیر ویژه، بردارهای ویژه مناسب را شناسایی کرد و تعداد خوشه‌ها را تخمین زد. این حالت ایده‌آل زمانی است که داده‌ها به وضوح از یکدیگر جدا باشند.  

اما در حالت واقعی، که خوشه‌ها کاملاً از یکدیگر جدا نیستند، همه بردارهای ویژه ممکن است مقادیر بزرگی داشته باشند. این موضوع باعث می‌شود که بردارهای ویژه با مقادیر ویژه بزرگ‌تر لزوماً مرتبط‌ترین یا مهم‌ترین بردارها نباشند. بنابراین، در چنین حالتی، صرفاً بررسی مقادیر ویژه کافی نیست و باید روش دیگری برای انتخاب بردارهای ویژه مناسب به کار گرفته شود.  

برای داده‌های واقعی، لازم است ارتباط هر بردار ویژه بر اساس ساختار و توزیع آن ارزیابی شود. روش پیشنهادی برای این منظور به بررسی توانایی هر بردار ویژه در جدا کردن داده‌ها به خوشه‌های مختلف می‌پردازد. این روش بدون هیچ فرضی درباره توزیع مجموعه داده اصلی \( D \)، ساختار بردارهای ویژه را مورد استفاده قرار می‌دهد تا ارتباط آن‌ها با خوشه‌بندی داده‌ها مشخص شود.  
\section{استفاده از توزیع گوسی}

برای ارزیابی میزان ارتباط یک بردار ویژه (\(e_k\))، از احتمالی به نام \(R_{ek}\) استفاده می‌شود که مقداری بین 0 و 1 دارد. اگر \(R_{ek} = 1\)، به این معناست که بردار ویژه بسیار مرتبط است، و اگر \(R_{ek} = 0\)، نشان می‌دهد که بردار ویژه غیرمرتبط است.

فرض می‌شود که عناصر یک بردار ویژه (\(e_k\)) می‌توانند یکی از دو نوع توزیع زیر را داشته باشند:
\begin{itemize}
	\item \textbf{توزیع تک‌مدی (Unimodal):} زمانی که بردار ویژه غیرمرتبط است.
	\item \textbf{توزیع چندمدی (Multimodal):} زمانی که بردار ویژه مرتبط است.
\end{itemize}

احتمال توزیع \(e_k\) بر اساس این دو حالت به صورت زیر تعریف می‌شود:
\[
p(e_k \mid \theta) = (1 - R_{ek})  p(e_{kn} \mid \theta_{e_{kn}}^1) + R_{ek}  p(e_{kn} \mid \theta_{e_{kn}}^2),
\]
که در آن:
\begin{itemize}
	\item \(\theta_{e_{kn}}^1\): پارامترهای توزیع تک‌مدی.
	\item \(\theta_{e_{kn}}^2\): پارامترهای توزیع چندمدی.
	\item \(R_{ek}\): وزن و احتمال مرتبط بودن \(e_k\).
\end{itemize}

برای توزیع تک‌مدی، فرض می‌شود که از یک توزیع نرمال (\(\mathcal{N}\)) با میانگین \(\mu_{k1}\) و واریانس \(\sigma_{k1}\) پیروی می‌کند:
\[
p(e_k \mid \theta_{e_{kn}}^1) = \mathcal{N}(e_k \mid \mu_{k1}, \sigma_{k1}).
\]

برای توزیع چندمدی، فرض می‌شود که از ترکیب دو توزیع نرمال با میانگین‌ها و واریانس‌های متفاوت تشکیل شده است:
\[
p(e_k \mid \theta_{e_{kn}}^2) = w_k \mathcal{N}(e_k \mid \mu_{k2}, \sigma_{k2}) + (1 - w_k) \mathcal{N}(e_k \mid \mu_{k3}, \sigma_{k3}),
\]
که در آن \(w_k\) وزن توزیع اول در این ترکیب است.

استفاده از ترکیب دو توزیع برای حالت چندمدی به این دلیل است که یک بردار ویژه مرتبط می‌تواند داده‌ها را به چند خوشه جدا کند و این ساختار چندمدی داده‌ها را بهتر نشان می‌دهد.

حال سوالی که ایجاد میشود این است که چرا از توزیع گوسی (نرمال) استفاده میکنیم؟

\textbf{1. تقارن:}  
توزیع نرمال یک توزیع متقارن است که به‌خوبی بسیاری از پدیده‌های طبیعی و داده‌های واقعی را مدل می‌کند. اگر داده‌ها حول یک مقدار مرکزی (میانگین) متقارن باشند، توزیع نرمال گزینه مناسبی است.

\textbf{2. سادگی در محاسبات:}  
توزیع نرمال دارای فرمول‌های بسته (Closed-Form) برای محاسبات میانگین، واریانس، و احتمال است که آن را برای مدل‌سازی و یادگیری ماشین بسیار ساده و مؤثر می‌سازد.

\textbf{3. پشتیبانی از بسیاری از الگوریتم‌ها:}  
بسیاری از روش‌های آماری و الگوریتم‌های یادگیری ماشین  بر پایه توزیع نرمال طراحی شده‌اند.

\vspace{0.5cm}

\textbf{قاعده حد مرکزی :}  
طبق این قاعده، مجموع یا میانگین تعداد زیادی متغیر تصادفی مستقل (حتی اگر از توزیع‌های غیر نرمال پیروی کنند) تقریباً از توزیع نرمال پیروی می‌کند. 


وقتی که \( e_{kn} \) بیش از دو خوشه تشکیل می‌دهد یا توزیع هر خوشه نرمال نیست، می‌توان از ترکیب دو توزیع گوسی (Gaussian) برای توصیف داده‌ها استفاده کرد. در این حالت:
\begin{enumerate}
	\item ترکیب دو توزیع گوسی (\( p(e_{kn} \mid \theta_{e_{kn}}^2) \)) نسبت به یک توزیع گوسی ساده بهتر به داده‌ها می‌نشیند.
	\item فرم ساده این ترکیب، تنها به تعداد محدودی از پارامترها برای توصیف \( p(e_{kn} \mid \theta_{e_{kn}}^2) \) نیاز دارد. این سادگی باعث می‌شود که یادگیری مدل حتی در صورت وجود داده‌های پراکنده نیز امکان‌پذیر باشد.
\end{enumerate}
\section{ML}

چند مشکل وجود دارد . ما از پیش نمی‌دانیم که آیا یک بردار ویژه مرتبط است یا غیرمرتبط.  
همچنین باید پارامترهای توزیع‌ها (مانند میانگین‌ها، واریانس‌ها و وزن‌های گوسی) را تخمین بزنیم.

برای توصیف توزیع \( e_{kn} \)، 8 پارامتر مورد نیاز است که به صورت زیر تعریف می‌شوند:
\[
\theta_{e_{kn}} = \{ R_{ek}, \mu_{k1}, \mu_{k2}, \mu_{k3}, \sigma_{k1}, \sigma_{k2}, \sigma_{k3}, w_k \}.
\]

در این وضعیت، حل مستقیم مسئله حداکثرسازی درستنمایی (ML) دشوار است، زیرا نمی‌دانیم \( e_{kn} \) به توزیع تک‌مدی یا چندمدی تعلق دارد (این اطلاعات "پنهان" است).  
الگوریتم EM با انجام دو مرحله به صورت متناوب به حل این مشکل کمک می‌کند:

\begin{itemize*}
	\item \textbf{گام E (Expectation):} حدس می‌زند که هر بخش از مدل (تک‌مدی یا چندمدی) چه مقدار مسئول توضیح \( e_{kn} \) است.
	\item \textbf{گام M (Maximization):} پارامترهای مدل (\( \theta \)) را بر اساس این حدس‌ها به‌روزرسانی می‌کند.
\end{itemize*}

\begin{itemize}
	\item ابتدا پارامترهای مؤلفه اول ترکیب (\( \theta_{e_{kn}}^1 \)) به صورت زیر براورد می‌شوند:
	\[
	\mu_{k1} = \frac{1}{N} \sum_{n=1}^N e_{kn}, \quad \sigma_{k1} = \frac{1}{N} \sum_{n=1}^N (e_{kn} - \mu_{k1})^2.
	\]
	\item سپس، 6 پارامتر باقی‌مانده به صورت تکراری با استفاده از روش بیشینه‌سازی انتظارات (Expectation-Maximization ) براورد می‌شوند.
	
	در روش حداکثرسازی انتظارات (EM)، به‌طور خاص، در گام \( E \)، احتمال پسین (Posterior-Probability) که هر مؤلفه ترکیب (Mixture-Component) مسئول \( e_{kn} \) است، به صورت زیر تخمین زده می‌شود:
	
	
	\[
	h_{kn}^1 = \frac{(1 - R_{ek}) \mathcal{N}(e_{kn} \mid \mu_{k1}, \sigma_{k1})}
	{(1 - R_{ek}) \mathcal{N}(e_{kn} \mid \mu_{k1}, \sigma_{k1}) + w_k R_{ek} \mathcal{N}(e_{kn} \mid \mu_{k2}, \sigma_{k2}) + (1 - w_k) R_{ek} \mathcal{N}(e_{kn} \mid \mu_{k3}, \sigma_{k3})},
	\]
	
	\[
	h_{kn}^2 = \frac{w_k R_{ek} \mathcal{N}(e_{kn} \mid \mu_{k2}, \sigma_{k2})}
	{(1 - R_{ek}) \mathcal{N}(e_{kn} \mid \mu_{k1}, \sigma_{k1}) + w_k R_{ek} \mathcal{N}(e_{kn} \mid \mu_{k2}, \sigma_{k2}) + (1 - w_k) R_{ek} \mathcal{N}(e_{kn} \mid \mu_{k3}, \sigma_{k3})},
	\]
	
	\[
	h_{kn}^3 = \frac{(1 - w_k) R_{ek} \mathcal{N}(e_{kn} \mid \mu_{k3}, \sigma_{k3})}
	{(1 - R_{ek}) \mathcal{N}(e_{kn} \mid \mu_{k1}, \sigma_{k1}) + w_k R_{ek} \mathcal{N}(e_{kn} \mid \mu_{k2}, \sigma_{k2}) + (1 - w_k) R_{ek} \mathcal{N}(e_{kn} \mid \mu_{k3}, \sigma_{k3})}.
	\]
	\begin{itemize}
		\item \( h_{kn}^1 \): مسئولیت جزء تک‌مدی.
		\item \( h_{kn}^2 \): مسئولیت گوسی اول در جزء چندمدی.
		\item \( h_{kn}^3 \): مسئولیت گوسی دوم در جزء چندمدی.
	\end{itemize}
	
	در گام \( M \) از روش حداکثرسازی انتظارات (EM)، شش پارامتر توزیع به‌صورت زیر بازتخمین زده می‌شوند:
	
	\[
	R_{ek}^{\text{new}} = 1 - \frac{1}{N} \sum_{n=1}^N h_{kn}^1, \quad w_k^{\text{new}} = \frac{1}{R_{ek}^{\text{new}} N} \sum_{n=1}^N h_{kn}^2,
	\]
	
	\[
	\mu_{k2}^{\text{new}} = \frac{\sum_{n=1}^N h_{kn}^2 e_{kn}}{\sum_{n=1}^N h_{kn}^2}, \quad \mu_{k3}^{\text{new}} = \frac{\sum_{n=1}^N h_{kn}^3 e_{kn}}{\sum_{n=1}^N h_{kn}^3},
	\]
	
	\[
	\sigma_{k2}^{\text{new}} = \frac{\sum_{n=1}^N h_{kn}^2 (e_{kn} - \mu_{k2}^{\text{new}})^2}{\sum_{n=1}^N h_{kn}^2}, \quad
	\sigma_{k3}^{\text{new}} = \frac{\sum_{n=1}^N h_{kn}^3 (e_{kn} - \mu_{k3}^{\text{new}})^2}{\sum_{n=1}^N h_{kn}^3}.
	\]
	
	الگوریتم EM ذاتاً یک روش جستجوی محلی (greedy) است و می‌تواند به مقداردهی اولیه پارامترها حساس باشد، به‌ویژه در داده‌های نویزی و پراکنده.  
	برای رفع این مشکل، مقدار اولیه \( R_{ek} \) به \( 0.5 \) تنظیم می‌شود و پنج پارامتر دیگر (\( \mu_{k2}, \mu_{k3}, \sigma_{k2}, \sigma_{k3}, w_k \)) به‌صورت تصادفی مقداردهی اولیه می‌شوند.  
	سپس، پاسخی که بالاترین مقدار \( p(e_{kn} \mid \theta_{e_{kn}}^2) \) را در میان مقداردهی‌های اولیه مختلف دارد، انتخاب می‌شود.
	
	\section{چرا هر فرمول به این شکل ساخته شده است؟}
	
	\subsection*{1. فرمول \( h_{kn}^1 \): مسئولیت جزء تک‌مدی}
	\[
	h_{kn}^1 = \frac{(1 - R_{ek}) \mathcal{N}(e_{kn} \mid \mu_{k1}, \sigma_{k1})}
	{(1 - R_{ek}) \mathcal{N}(e_{kn} \mid \mu_{k1}, \sigma_{k1}) + w_k R_{ek} \mathcal{N}(e_{kn} \mid \mu_{k2}, \sigma_{k2}) + (1 - w_k) R_{ek} \mathcal{N}(e_{kn} \mid \mu_{k3}, \sigma_{k3})}.
	\]
	این فرمول احتمال اینکه \( e_{kn} \) توسط جزء تک‌مدی (بردارهای ویژه غیرمرتبط) توضیح داده شود را محاسبه می‌کند.
	
	\begin{itemize}
		\item \textbf{صورت کسر:} 
		\((1 - R_{ek}) \mathcal{N}(e_{kn} \mid \mu_{k1}, \sigma_{k1})\)،
		که شامل احتمال غیرمرتبط بودن (\( 1 - R_{ek} \)) و احتمال توزیع گوسی تک‌مدی است.
		\item \textbf{مخرج کسر:} جمع احتمالات همه اجزا (تک‌مدی و چندمدی) برای نرمال‌سازی.
		\item \textbf{ایده اصلی:} 
		اگر \( e_{kn} \) با توزیع تک‌مدی مطابقت بیشتری داشته باشد، مقدار \( h_{kn}^1 \) بزرگ‌تر خواهد بود.
	\end{itemize}
	\( h_{kn}^3 \)و\( h_{kn}^2 \) نیز به طور مشابه تعریف میشوند
	
	\subsection*{2. فرمول \( R_{ek}^{\text{new}} \): مرتبط بودن بردار ویژه}
	\[
	R_{ek}^{\text{new}} = 1 - \frac{1}{N} \sum_{n=1}^N h_{kn}^1.
	\]
	این فرمول نشان می‌دهد که بردار ویژه \( e_k \) چقدر احتمال دارد مرتبط باشد.
	
	\begin{itemize}
		\item \textbf{میانگین \( h_{kn}^1 \):} احتمال کلی غیرمرتبط بودن.
		\item \textbf{کم کردن از 1:} 
		این کار احتمال مرتبط بودن را به دست می‌آورد.
		\item \textbf{ایده اصلی:} 
		اگر اکثر \( e_{kn} \) توسط جزء تک‌مدی توضیح داده شوند، مقدار \( R_{ek} \) کوچک خواهد شد.
	\end{itemize}
	

	
	\subsection*{3. فرمول \( \mu_{k2}^{\text{new}} \): میانگین گوسی اول در جزء چندمدی}
	\[
	\mu_{k2}^{\text{new}} = \frac{\sum_{n=1}^N h_{kn}^2 e_{kn}}{\sum_{n=1}^N h_{kn}^2}.
	\]
	این فرمول میانگین وزنی \( e_{kn} \) را برای گوسی اول محاسبه می‌کند.
	
	\begin{itemize}
		\item \textbf{صورت کسر:} 
		\( e_{kn} \) با احتمال تعلق به گوسی اول (\( h_{kn}^2 \)) وزن‌دهی می‌شود.
		\item \textbf{مخرج کسر:} 
		مجموع وزن‌ها برای نرمال‌سازی.
		\item \textbf{ایده اصلی:} 
		داده‌هایی که احتمال تعلق بیشتری به گوسی اول دارند، تأثیر بیشتری بر میانگین خواهند داشت.
	\end{itemize}
	

	
	\subsection*{4. فرمول \( \sigma_{k2}^{\text{new}} \): واریانس گوسی اول در جزء چندمدی}
	\[
	\sigma_{k2}^{\text{new}} = \frac{\sum_{n=1}^N h_{kn}^2 (e_{kn} - \mu_{k2}^{\text{new}})^2}{\sum_{n=1}^N h_{kn}^2}.
	\]
	این فرمول واریانس وزنی \( e_{kn} \) را برای گوسی اول محاسبه می‌کند.
	
	\begin{itemize}
		\item \textbf{صورت کسر:} 
		انحراف مربع \( e_{kn} \) از میانگین (\( \mu_{k2} \))، با احتمال تعلق (\( h_{kn}^2 \)) وزن‌دهی می‌شود.
		\item \textbf{مخرج کسر:} 
		مجموع وزن‌ها برای نرمال‌سازی.
		\item \textbf{ایده اصلی:} 
		واریانس منعکس‌کننده پراکندگی داده‌های مرتبط با گوسی اول است.
	\end{itemize}
	
	
	
	\subsection*{5. فرمول \( w_k^{\text{new}} \): وزن گوسی اول}
	\[
	w_k^{\text{new}} = \frac{1}{R_{ek}^{\text{new}} N} \sum_{n=1}^N h_{kn}^2.
	\]
	این فرمول سهم گوسی اول را در جزء چندمدی مشخص می‌کند.
	
	\begin{itemize}
		\item \textbf{صورت کسر:} 
		مجموع مسئولیت‌های گوسی اول (\( h_{kn}^2 \)).
		\item \textbf{مخرج کسر:} 
		\( R_{ek} N \) برای نرمال‌سازی سهم گوسی اول.
		\item \textbf{ایده اصلی:} 
		اگر گوسی اول بیشتر داده‌ها را توضیح دهد، وزن آن افزایش می‌یابد.
	\end{itemize}
	
	
	
	\section{وزن‌ها و اهمیت آن‌ها در مدل}
	
	وزن‌ها (\(w_k\)) مقادیر احتمالی هستند که اهمیت نسبی دو گوسی در جزء چندمدی را توصیف می‌کنند.  
	در اینجا:
	\begin{itemize}
		\item \(w_k\): وزن گوسی اول است که نشان می‌دهد چه مقدار از داده‌ها توسط گوسی اول توضیح داده می‌شود.
		\item \(1 - w_k\): وزن گوسی دوم است که سهم آن را در توضیح داده‌ها نشان می‌دهد.
	\end{itemize}
	
	وزن‌ها در بازه \(0\) تا \(1\) قرار دارند و تضمین می‌شود که مجموع آن‌ها برابر با \(1\) باشد:
	\[
	w_k + (1 - w_k) = 1.
	\]
	
	
	
	در جزء چندمدی، توزیع گوسی چندمدی ترکیبی از دو گوسی است:
	\[
	p(e_{kn} \mid \theta_{\text{multimodal}}) = w_k \mathcal{N}(e_{kn} \mid \mu_{k2}, \sigma_{k2}) + (1 - w_k) \mathcal{N}(e_{kn} \mid \mu_{k3}, \sigma_{k3}).
	\]
	در این فرمول:
	\begin{itemize}
		\item \(\mathcal{N}(e_{kn} \mid \mu_{k2}, \sigma_{k2})\): گوسی اول که یکی از خوشه‌های داده‌های مرتبط با بردار ویژه را توضیح می‌دهد.
		\item \(\mathcal{N}(e_{kn} \mid \mu_{k3}, \sigma_{k3})\): گوسی دوم که خوشه دیگری از داده‌ها را توضیح می‌دهد.
	\end{itemize}
	
	
	وزن‌ها به مدل انعطاف‌پذیری می‌دهند تا بتواند به خوبی توزیع داده‌ها را مدیریت کند:
	\begin{itemize}
		\item اگر بیشتر داده‌های مرتبط با بردار ویژه به گوسی اول تعلق داشته باشند، مقدار \(w_k\) به \(1\) نزدیک خواهد شد.
		\item اگر داده‌ها به طور مساوی بین دو گوسی تقسیم شده باشند، مقدار \(w_k\) نزدیک به \(0.5\) خواهد بود.
	\end{itemize}
	
	این انعطاف به مدل اجازه می‌دهد تا داده‌هایی با حالت‌های چندگانه (multimodal) را به خوبی مدیریت کرده و سهم هر گوسی در جزء چندمدی را به طور دقیق تنظیم کند.
	
\end{itemize}


اگرچه الگوریتم یادگیری ارتباط ما بر اساس تخمین توزیع عناصر هر بردار ویژه است، ما تنها به یادگیری احتمال اینکه توزیع تک‌مدی یا چندمدی باشد علاقه‌مندیم، که این موضوع با مقدار \( R_{ek} \) نشان داده می‌شود.  
به عبارت دیگر، از میان ۸ پارامتر آزاد توزیع بردار ویژه  تنها پارامتر \( R_{ek} \) است که ما به دنبال آن هستیم. به همین دلیل است که الگوریتم ما عملکرد خوبی دارد.
\begin{figure}[h!]
	\centering

	\includegraphics[width=0.9\textwidth]{a.png}\par\vspace{1cm}
 \caption{
	بردارهای ویژه مصنوعی و معیار ارتباط تخمینی \( R_{ek} \). عناصر هر بردار ویژه از سه خوشه که به صورت یکنواخت توزیع شده‌اند تشکیل شده است. میانگین سه خوشه به ترتیب \( m1, 0, -m1 \) هستند. بدیهی است که هر چه مقدار \( m1 \) بزرگ‌تر باشد، سه خوشه در توزیع بردار ویژه متمایزتر می‌شوند و بردار ویژه مرتبط‌تر خواهد بود. 
	(a) \( m1 = 0.0, R_{ek} = 0.2992 \). 
	(b) \( m1 = 0.2, R_{ek} = 0.4479 \). 
	(c) \( m1 = 0.4, R_{ek} = 0.6913 \). 
	(d) \( m1 = 0.6, R_{ek} = 0.7409 \). 
	(e) \( m1 = 0.8, R_{ek} = 0.8302 \). 
	(f) \( m1 = 1.0, R_{ek} = 1.0 \).
}
\label{fig:synthetic_eigenvectors}
\end{figure}

حتی زمانی که بیش از دو خوشه وجود دارد یا توزیع هر خوشه گوسی نیست، الگوریتم ما همچنان می‌تواند به درستی مرتبط بودن یا اطلاع‌رسان بودن یک بردار ویژه را نشان دهد.  
ا 
به طور خاص، شکل بالا نشان می‌دهد که زمانی که توزیع عناصر بردار ویژه متعلق به خوشه‌های مختلف یکنواخت است و بیش از دو خوشه وجود دارد، مقادیر \( R_{ek} \) که با استفاده از الگوریتم ما تخمین زده شده‌اند، همچنان می‌توانند به درستی منعکس‌کننده مرتبط بودن یک بردار ویژه باشند.  


\section{حذف و وزن‌دهی بردارهای ویژه}
\( R_{ek} \) تخمینی، اندازه‌گیری پیوسته‌ای از میزان مرتبط بودن \( e_k \) ارائه می‌دهد.  
از آنجا که برای کاهش ابعاد به یک "تصمیم قطعی" نیاز است، به سادگی \( e_k \) را از میان \( K_m \) بردار ویژه نامزد حذف می‌کنیم اگر:
\[
R_{ek} < 0.5. 
\]

بردارهای ویژه مرتبط باقی‌مانده سپس با استفاده از \( R_{ek} \) وزن‌دهی می‌شوند.  
این کار یک مجموعه داده جدید به صورت زیر تولید می‌کند:
\[
D_r = \{y_1, \dots, y_n, \dots, y_N\}, 
\]
که در آن \( y_n \) یک بردار ویژگی با ابعاد \( K_r \) است که برابر با تعداد بردارهای ویژه مرتبط انتخاب‌شده است.  
ما توزیع \( D_r \) را با استفاده از یک مدل ترکیبی گوسی (GMM) برای خوشه‌بندی داده‌ها مدل‌سازی می‌کنیم.  
سپس از BIC برای انتخاب تعداد بهینه مؤلفه‌های گوسی استفاده می‌کنیم که با تعداد بهینه خوشه‌ها \( K_o \) مطابقت دارد.  
هر بردار ویژگی در مجموعه داده آموزشی سپس به عنوان یکی از \( K_o \) خوشه‌ها با استفاده از GMM آموزش‌دیده با \( K_o \) مؤلفه‌های گوسی برچسب‌گذاری می‌شود.  


\section{الگوریتم پیشنهادی خوشه‌بندی طیفی بر اساس انتخاب بردار ویژه مرتبط}

\textbf{ورودی:} یک مجموعه از \( N \) نقطه داده/الگوی ورودی \( D \).  

\textbf{خروجی:} تعداد بهینه خوشه‌ها \( K_o \) و \( D \) که به \( K_o \) زیرمجموعه تقسیم شده است.  

\begin{enumerate}
	\item تشکیل ماتریس وابستگی \( A = \{A_{ij}\} \).
	\item ساخت ماتریس وابستگی نرمال‌شده \( \tilde{A} \).
	\item یافتن \( K_m \) بزرگ‌ترین بردارهای ویژه \( \tilde{A} \)، یعنی \( e_1, e_2, \dots, e_{K_m} \)، و تشکیل \( D_e \).
	\item تخمین میزان مرتبط بودن هر بردار ویژه \( R_{ek} \) با استفاده از الگوریتم انتخاب بردار ویژه پیشنهادی.
	\item حذف بردار ویژه \( e_k \) اگر \( R_{ek} < 0.5 \).
	\item وزن‌دهی بردارهای ویژه مرتبط با استفاده از \( R_{ek} \) و تشکیل یک مجموعه داده جدید \( D_r \).
	\item مدل‌سازی توزیع \( D_r \) با استفاده از مدل ترکیبی گوسی (GMM) و تخمین تعداد مؤلفه‌های گوسی به عنوان \( K_o \) با استفاده از BIC.
	\item اختصاص نقطه داده اصلی \( f_i \) به خوشه \( j \) اگر و فقط اگر نقطه \( i \) در \( D_r \) به خوشه \( j \) اختصاص داده شده باشد.
\end{enumerate}

\chapter{مقایسه روش های خوشه بندی}

\section*{1. خوشه‌بندی (K-Means):}

\subsection*{مزایا:}
\begin{itemize}
	\item  الگوریتم \( k \)-میانگین بسیار ساده و محاسباتی سریع است، به خصوص برای داده‌های بزرگ با تعداد ویژگی کم.
	\item به راحتی می‌تواند برای مجموعه داده‌های بزرگ و با ابعاد زیاد استفاده شود.
	\item  بهترین عملکرد را زمانی دارد که داده‌ها به شکل خوشه‌های کروی و با اندازه مشابه هستند.
\end{itemize}

\subsection*{معایب:}
\begin{itemize}
	\item  تعداد خوشه‌ها باید از قبل مشخص شود که در بسیاری از مواقع این مقدار ناشناخته است.
	\item  اگر خوشه‌ها از شکل‌های پیچیده یا غیریکسان برخوردار باشند، \( k \)-میانگین نمی‌تواند مرزهای خوشه‌ها را به خوبی شناسایی کند.
	\item  مقادیر پرت یا نویز می‌توانند تأثیر زیادی بر موقعیت مراکز خوشه (\textit{centroids}) داشته باشند.
\end{itemize}

\section*{2. خوشه‌بندی طیفی (Spectral-Clustering):}

\subsection*{مزایا:}
\begin{itemize}
	\item  خوشه‌بندی طیفی می‌تواند خوشه‌هایی با شکل‌های غیرخطی (مانند حلقه‌ها یا مارپیچ‌ها) را شناسایی کند.
	\item از اطلاعات مربوط به شباهت داده‌ها (مانند یک ماتریس شباهت یا گراف) استفاده می‌کند، بنابراین برای داده‌هایی که ساختار پیچیده دارند بسیار موثر است.
	\item  برخلاف \( k \)-میانگین، خوشه‌بندی طیفی نیازی به فرض کروی بودن خوشه‌ها ندارد.
\end{itemize}

\subsection*{معایب:}
\begin{itemize}
	\item خوشه‌بندی طیفی برای مجموعه داده‌های بزرگ (با تعداد نمونه‌های زیاد) بسیار کند است، زیرا محاسبه مقادیر ویژه ماتریس شباهت هزینه‌بر است.
	\item  انتخاب مناسب پارامترهایی مانند سیگما (\( \sigma \)) در ماتریس شباهت تأثیر زیادی بر نتایج دارد.
	\item  همانند \( k \)-میانگین، تعداد خوشه‌ها باید از قبل مشخص شود.
\end{itemize}

\chapter{تحلیل داده ها}
در این بخش سه پیاده سازی متفاوت قرار داده ایم ؛ روش k-means و روش  خوشه بندی طیفی که یکبار آنرا با کتابخانه آماده منتشر کرده ایم و بار دیگر به صورت دستی و مرحله به مرحله . در ابتدا هر الگوریتم را با تعداد 7 خوشه مقایسه میکنیم سپس عملکردشان در تعداد کمتری از خوشه یعنی 4 عدد را بررسی میکنیم .

در روش  (GMM)، از معیار اطلاعات بیزی (BIC) استفاده می‌شود تا تعداد بهینه خوشه‌ها را از میان بازه‌ای مشخص انتخاب کند (در کد، بازه 2 تا 10). فرآیند به این صورت است که مدل GMM برای مقادیر مختلف \( k \) (تعداد خوشه‌ها) آموزش داده می‌شود و مقدار BIC برای هر مقدار \( k \) محاسبه می‌گردد:

\[
\text{BIC} = -2 \cdot \log(\text{Likelihood}) + \text{Penalty\_Term}
\]

این معیار، مدلی را ترجیح می‌دهد که بیشترین \textit{Likelihood} را داشته باشد، اما پیچیدگی بیشتر (تعداد پارامترهای بالاتر) را جریمه می‌کند. تعداد خوشه‌ای که کمترین مقدار BIC را داشته باشد، به عنوان تعداد بهینه انتخاب می‌شود.

در داده‌های تولید‌شده ، به دلیل مقدار زیاد پراکندگی خوشه‌ها ،خوشه‌ها همپوشانی زیادی دارند. برای مدل‌سازی بهتر این داده‌های پیچیده، GMM داده‌ها را به 7 خوشه تقسیم کرد تا توزیع داده‌ها را بهتر مدل کند. به طور خاص، در GMM فرض بر این است که داده‌ها از توزیع نرمال چندمتغیره پیروی می‌کنند، و به دلیل همپوشانی داده‌ها، مدل برای پوشش این ساختار به تعداد بیشتری خوشه نیاز داشت.

در مقابل، زمانی که تعداد خوشه‌ها به صورت دستی ثابت می‌شود (\( k =4 \))، GMM دیگر تعداد بهینه را پیدا نمی‌کند و صرفاً داده‌ها را در چهار خوشه توزیع می‌کند. این کار باعث می‌شود که مدل ساختار داده‌ها را با تعداد خوشه‌های مشخص‌شده تطبیق دهد، حتی اگر داده‌ها به تعداد بیشتری خوشه نیاز داشته باشند. بنابراین، در این روش تعداد خوشه‌ها به صورت صریح توسط خودمان مشخص می‌شود، در حالی که در روش قبلی، تعداد خوشه‌ها به صورت خودکار و پویا تعیین می‌شد.

از طرفی در خوشه بندی طیفی با کتابخانه ابتدا، الگوریتم مقادیر ویژه و بردارهای ویژه ماتریس لاپلاسین را محاسبه می‌کند. سپس، تنها \(k\) بردار ویژه اول (بردارهایی که مقادیر ویژه کوچک‌تری دارند) انتخاب می‌شوند. این بردارها داده‌ها را به فضایی جدید منتقل می‌کنند که در آن خوشه‌ها بهتر از هم جدا شده‌اند.

پس از انتقال داده‌ها به این فضای جدید، الگوریتم k-means یا یک روش مشابه برای خوشه‌بندی اعمال می‌شود. در این مرحله، تعداد خوشه‌ها (\(k\)) به صورت مستقیم به الگوریتم k-means پاس داده می‌شود تا داده‌ها به همان تعداد خوشه تقسیم شوند. این فرآیند تضمین می‌کند که تعداد نهایی خوشه‌ها دقیقاً برابر مقدار تعیین‌شده توسط کاربر باشد.
وقتی ما از خوشه‌بندی طیفی با GMM استفاده می‌کنیم و تعداد خوشه‌ها را ثابت می‌کنیم، الگوریتم داده‌ها را در فضایی که با بردارهای ویژه ساخته شده است، مدل می‌کند. این مدل بر اساس توزیع نرمال چندمتغیره عمل کرده و می‌تواند خوشه‌های پیچیده یا همپوشانی را بهتر مدیریت کند. اما وقتی از خوشه‌بندی طیفی کتابخانه‌ای مثل Scikit-learn استفاده می‌کنیم، تعداد خوشه‌ها از قبل مشخص می‌شود و الگوریتم \(k\)-Means داده‌ها را در خوشه‌های کروی و جدا از هم قرار می‌دهد. این روش ساده‌تر و سریع‌تر است، اما در داده‌های پیچیده یا با خوشه‌های همپوشانی عملکرد ضعیف‌تری دارد.



\begin{figure}[h!]
	\centering
	\includegraphics[width=0.9\textwidth]{b.png}
	\caption{داده‌های خام مورد بررسی}
\end{figure}

\begin{figure}[h!]
	\centering
	\includegraphics[width=0.7\textwidth]{c.png}
	\caption{خوشه‌بندی طیفی سفارشی با 7 خوشه}
\end{figure}

\begin{figure}[h!]
	\centering
	\includegraphics[width=0.7\textwidth]{z.png}
	\caption{خوشه‌بندی k-means با 7 خوشه}
\end{figure}

\begin{figure}[h!]
	\centering
	\includegraphics[width=0.7\textwidth]{v.png}
	\caption{خوشه‌بندی طیفی با استفاده از کتابخانه با 7 خوشه}
\end{figure}

\begin{figure}[h!]
	\centering
	\includegraphics[width=0.7\textwidth]{sa.png}
	\caption{خوشه‌بندی طیفی سفارشی با 4 خوشه}
\end{figure}

\begin{figure}[h!]
	\centering
	\includegraphics[width=0.7\textwidth]{u.png}
	\caption{خوشه‌بندی k-means با 4 خوشه}
\end{figure}

\begin{figure}[h!]
	\centering
	\includegraphics[width=0.7\textwidth]{o.png}
	\caption{خوشه‌بندی طیفی با استفاده از کتابخانه با 4 خوشه}
\end{figure}

\begin{figure}[h!]
	\centering
	\includegraphics[width=0.7\textwidth]{q.png}
	\caption{ماتریس شباهت قبل از سورت شدن در روش خوشه‌بندی طیفی سفارشی}
\end{figure}

\begin{figure}[h!]
	\centering
	\includegraphics[width=0.7\textwidth]{f.png}
	\caption{ماتریس شباهت بعد از سورت شدن در روش خوشه‌بندی طیفی سفارشی}
\end{figure}
\chapter{کاربردها}
خوشه‌بندی طیفی به دلیل توانایی شناسایی الگوهای پیچیده و غیرخطی، کاربردهای گسترده و جذابی در علوم و صنعت دارد. در پردازش تصویر، برای تفکیک بخش‌های مختلف تصاویر و ردیابی اشیاء به‌کار می‌رود. در تحلیل شبکه‌های اجتماعی، این روش گروه‌های اجتماعی و افراد تأثیرگذار را آشکار می‌کند. در پردازش زبان طبیعی، ابزار قدرتمندی برای دسته‌بندی اسناد و کشف ارتباطات معنایی میان واژه‌هاست. در زیست‌محاسبات، تحلیل داده‌های ژنومی و مدل‌سازی شبکه‌های زیستی را تسهیل می‌کند. در تحلیل سیگنال و سری‌های زمانی، خوشه‌بندی داده‌ها و تشخیص ناهنجاری‌ها را ممکن می‌سازد.

این روش در یادگیری ماشین، به‌عنوان ابزاری برای کاهش ابعاد داده‌های پیچیده و پیش‌پردازش کارآمد مطرح است. در سیستم‌های توصیه‌گر، شناسایی علایق کاربران و ارائه پیشنهادات شخصی‌سازی‌شده را ارتقا می‌دهد. در شبکه‌های حمل‌ونقل، مدیریت ترافیک و طراحی بهینه مسیرها را ممکن می‌سازد. در حوزه امنیت، رفتارهای غیرمعمول و حملات سایبری را شناسایی کرده و در مهندسی، برای تحلیل داده‌های حسگرها و بهینه‌سازی مصرف انرژی به‌کار گرفته می‌شود.

این تنوع کاربرد، خوشه‌بندی طیفی را به ابزاری قدرتمند و انعطاف‌پذیر برای حل مسائل پیچیده در دنیای امروز تبدیل کرده است.
در این پروژه با استفاده از خوشه بندی طیفی به تقسیم بندی تصاویر و شناسایی اجزای داخل تصاویر میپردازیم که نتایج آن آورده شده است.
در گاربرد اول به نوعی اجزای تصویرمان را آبجکتیفای میکنیم .این فرآیند خوشه‌بندی طیفی را برای تقسیم‌بندی تصویر سکه‌ها به نواحی مختلف به‌کار می‌برد. ابتدا تصویر صاف و کوچک‌تر می‌شود، سپس به گرافی تبدیل می‌گردد که شباهت بین پیکسل‌ها را بر اساس شدت روشنایی مدل‌سازی می‌کند. الگوریتم خوشه‌بندی طیفی با سه روش مختلف برای تخصیص برچسب‌ها اجرا شده و تصویر به نواحی خوشه‌بندی‌شده تقسیم می‌شود.
در کاربرد بعدی تصویری مصنوعی شامل چند مستطیل و بیضی تولید می‌شود و نویز به آن افزوده می‌گردد. سپس، تصویر به یک گراف تبدیل شده و خوشه‌بندی طیفی برای تقسیم‌بندی آن به ۶ ناحیه اعمال می‌شود. در نهایت، تصویر اصلی و نسخه خوشه‌بندی‌شده در کنار هم نمایش داده می‌شوند.
	\begin{figure}[h!]
		\centering
		\includegraphics[width=0.7\textwidth]{x.png}
	\end{figure}
	
	\begin{figure}[h!]
		\centering
		\includegraphics[width=0.7\textwidth]{y.png}
	\end{figure}
	
	\begin{figure}[h!]
		\centering
		\includegraphics[width=0.7\textwidth]{k.png}
		\caption{کاربرد اول}
	\end{figure}
	
	\begin{figure}[h!]
		\centering
		\includegraphics[width=0.9\textwidth]{p.png}
		\caption{کاربرد دوم}
	\end{figure}
\chapter{بحث و نتیجه‌گیری}
\LTR
\begin{table}[h!]
	\centering
	\begin{tabular}{@{}lcc@{}}
		\toprule
		\textbf{الگوریتم}               & \textbf{ k=7} & \textbf{  k=4} \\ \midrule
		Spectral Clustering (GMM)       &         40.0      &          52.0     &                               \\ 
		Spectral Clustering (Library)   &           35.0         &          49.0 &                               \\ 
		K-Means                         &           35.0          &         52.0 &                               \\ \bottomrule
	\end{tabular}
	
.
	\RTL
	\caption{مقایسه ضرایب سیلوئت برای سه الگوریتم خوشه‌بندی}
	\label{tab:silhouette_comparison}
	جدول فوق نشان میدهد که با هر دو انتخاب خوشه خوشه بندی سفارشی ما عملکرد بهتر یا هم ترازی از خود نشان میدهد. روش کتابخانه ای نیز از دقت خوبی برخوردار بوده و در عین حال در پیاده سازی سادگی و راحتی بیشتری دارد. در نهایت روش کلاسیک k-means در تعداد خوشه کم به خوبی عمل میکند اما با افزایش خوشه دقتش افت میکند.
\end{table}
\RTL
\chapter{پیوست ها}
	\begin{figure}[h!]
	\centering
	\includegraphics[width=0.5\textwidth]{1.7.1.png}


	\includegraphics[width=0.5\textwidth]{1.7.2.png}

	
	\includegraphics[width=0.5\textwidth]{1.7.3.png}
	\caption{خوشه بندی با مدل گوسی و 7 خوشه}
\end{figure}
\begin{figure}[h!]
	\centering
	\includegraphics[width=0.5\textwidth]{1.4.1.png}


	\includegraphics[width=0.5\textwidth]{1.4.2.png}

	\centering
	\includegraphics[width=0.5\textwidth]{1.4.3.png}
	\caption{خوشه بندی با مدل گوسی و 4 خوشه}
\end{figure}\begin{figure}[h!]
\centering
\includegraphics[width=0.7\textwidth]{2.4.png}
\caption{خوشه بندی با k-means}
\end{figure}\begin{figure}[h!]
\centering
\includegraphics[width=0.7\textwidth]{3.4.png}
\caption{خوشه بندی طیفی با کتابخانه}
\end{figure}\begin{figure}[h!]
\centering
\includegraphics[width=0.7\textwidth]{coin.png}
\caption{آبجکتیفای کردن سکه ها}
\end{figure}\begin{figure}[h!]
\centering
\includegraphics[width=1\textwidth]{rec.png}
\caption{تقسیم بندی تصویر}
\end{figure}
	\LTR
	\begin{flushleft} 
		\begin{thebibliography}{9}
			
			\bibitem{xiang2007spectral}
			Tao Xiang and Shaogang Gong. \textit{Spectral clustering with eigenvector selection.} Pattern Recognition, vol. 41, no. 3, pp. 1012--1029, 2007.
			\bibitem{ning2010incremental}
			Huazhong Ning, Wei Xu, Yun Chi, Yihong Gong, and Thomas S. Huang.
			\textit{Incremental spectral clustering by efficiently updating the eigen-system.}
			Pattern Recognition, vol. 43, no. 1, pp. 113--127, 2010.
			
			
			\bibitem{luxburg2007tutorial}
			Ulrike von Luxburg.
			\textit{A Tutorial on Spectral Clustering.} 
			Technical Report No. TR-149, Max Planck Institute for Biological Cybernetics, March 2007.
			
			\bibitem{ng2002spectral}
			Andrew Y. Ng, Michael I. Jordan, and Yair Weiss.
			\textit{On Spectral Clustering: Analysis and an Algorithm.} 
			Advances in Neural Information Processing Systems (NIPS), vol. 14, pp. 849--856, 2002.
			
			
			\bibitem{grosse_mixture}
			Roger Grosse and Nitish Srivastava.
			\textit{Lecture 16: Mixture Models.} 
			Available online at: \url{https://www.cs.toronto.edu/~rgrosse/csc321/mixture_models.pdf}.
			
			
			\bibitem{ghasempour_clustering}
			Kamyar Ghasemipour.
			\textit{Clustering, K-Means, EM Tutorial.} 
			Available online at: \url{https://www.cs.toronto.edu/~jlucas/teaching/csc411/lectures/tut8_handout.pdf}.
			
			\bibitem{laplacian_matrix}
			Wikipedia contributors.
			\textit{Laplacian matrix.} 
			Available online at: \url{https://en.wikipedia.org/wiki/Laplacian_matrix}, accessed January 2025.
			\bibitem{geeksforgeeks_spectral}
			GeeksforGeeks.
			\textit{ML - Spectral Clustering.} 
			Available online at: \url{https://www.geeksforgeeks.org/ml-spectral-clustering/}, accessed January 2025.
			\bibitem{lavan_clustering}
			Lavan Ertebat.
			\textit{Clustering Algorithms.} 
			Available online at: \url{https://lavanertebat.com/clustering-algorithms/}, accessed January 2025.
			\bibitem{scikit_spectral}
			Scikit-learn.
			\textit{Spectral Clustering.} 
			Available online at: \url{https://scikit-learn.org/1.5/modules/generated/sklearn.cluster.SpectralClustering.html}, accessed January 2025.
			\bibitem{deepai_affinity}
			DeepAI.
			\textit{Affinity Matrix.} 
			Available online at: \url{https://deepai.org/machine-learning-glossary-and-terms/affinity-matrix}, accessed January 2025.
		\end{thebibliography}
	\end{flushleft} 
	
	
\end{document}
